\SourcePage{764}{767}
\section{Neutron scattering}

In a number of collision-theory problems, it is necessary to determine how the scattering process is affected by the motion of the scattering centers themselves. Under certain conditions, such problems can be treated by a distinctive perturbation theory developed by Fermi in 1936, even though perturbation theory need not apply to scattering by each center in isolation. One problem of this kind is the scattering of slow neutrons by a system of atoms, such as a molecule.

For definiteness, this is the particular problem discussed below.

Electrons scatter neutrons to a negligible extent, so in practice all the scattering occurs at the nuclei\footnote{It is also assumed that the molecule has no magnetic moment. Otherwise there is an additional, specific scattering effect due to the interaction between the magnetic moments of the molecule and the neutron.}. Suppose that the scattering amplitude of an individual nucleus is small compared with the interatomic distances. Then the amplitude of the wave

\SourcePage{765}{768}
scattered by any one nucleus in the molecule is already small at the positions of the other nuclei. Under these conditions, the scattering amplitude of the molecule reduces to the sum of the amplitudes for scattering by the individual nuclei.

Perturbation theory is generally inapplicable to a collision between a neutron and a nucleus: although nuclear forces have a short range, they are very strong within that range. What matters, however, is that the scattering amplitude of a slow neutron, whose wavelength is large compared with the nuclear dimensions, is a constant independent of velocity. Let $f_a$ be the amplitude for scattering by the $a$th nucleus; then $|f_a|^2do$ is the differential cross section for elastic scattering of the neutron by a free nucleus, in their center-of-mass system.

A constant amplitude can be obtained formally from perturbation theory if the neutron--nucleus interaction is described by the ``point'' potential energy
\begin{equation}
 U(\mathbf r)=-\frac{2\pi\hbar^2}{M}f\delta(\mathbf r),
 \tag{151.1}
\end{equation}
where $M$ is the reduced mass of the neutron and nucleus. When this expression is substituted into the Born formula (126.4), the delta-function turns the integral into a constant independent of $\mathbf q$. The ``field'' $U(\mathbf r)$ so defined is called a \emph{pseudopotential}. Its introduction is possible precisely because $f$ is constant.

At an arbitrary neutron energy, the scattering amplitude generally depends separately on the initial and final momenta $\mathbf p$ and $\mathbf p'$, rather than only on their difference $\mathbf q$; an amplitude calculated in the Born approximation, by contrast, can depend only on $\mathbf q$\footnote{It should also be stressed that although the pseudopotential gives the correct scattering amplitude when perturbation theory is applied formally, this by no means implies that perturbation theory genuinely applies to such a field. On the contrary, for a potential well whose depth $U_0$ tends to infinity according to $U_0a^3=\mathrm{const}$, where the well radius $a$ tends to zero, conditions (126.1) and (126.2) are certainly not satisfied.}.

If the scattering nucleus undergoes a specified motion, for example vibration within a molecule, averaging the interaction (151.1) over this motion ``smears'' it across a region whose dimensions are generally large compared with the scattering amplitude $f$. The smeared interaction satisfies condition (126.1) for applicability of the Born approximation.

\SourcePage{766}{769}
We shall therefore describe the neutron's interaction with the molecule by the pseudopotential
\begin{equation}
 U(\mathbf r)=-2\pi\hbar^2\sum_a\frac{f_a}{M_a}\delta(\mathbf r-\mathbf R_a),
 \tag{151.2}
\end{equation}
where the sum is over all nuclei in the molecule, $\mathbf R_a$ are their radius vectors, and $\mathbf r$ is the neutron radius vector. Substitution into the perturbation-theory formula (148.3), with the reduced mass $M_\text{mol}$ of the molecule and neutron used for $m$, gives the following neutron--molecule scattering cross section in their center-of-mass system:
\begin{equation}
 d\sigma_n=M_\text{mol}^2\frac{p'}p
 \left|\sum_a\frac{f_a}{M_a}\langle n|e^{-i\mathbf q\mathbf R_a}|0\rangle\right|^2do.
 \tag{151.3}
\end{equation}
The matrix elements are taken between wave functions of stationary states of the nuclear motion, with energies $E_0$ and $E_n$. The momenta $p$ and $p'$ are related by energy conservation:
\[
 \frac{p^2-p'^2}{2M_\text{mol}}=E_n-E_0.
\]

Formula (151.3) describes an inelastic collision accompanied by a specified change in the state of nuclear motion in the molecule, the transition $0\longrightarrow n$. It solves the problem posed: from the assumed known neutron-scattering amplitudes of the free nuclei, it determines the molecular scattering cross section while accounting for the nuclei's own motion and for interference between scattering at different nuclei.

If the nuclei have nonzero spin, one must additionally allow for the dependence of the amplitudes $f_a$ on the resultant spin of the scattering nucleus and neutron. This can be done as follows.

The resultant spin of a nucleus and neutron has either of the values $j_a=i_a\pm1/2$, where $i_a$ is the nuclear spin. Denote the corresponding scattering amplitudes by $f_a^+$ and $f_a^-$. We construct a spin operator whose eigenvalues at the specified values of $j_a$ are $f_a^+$ and $f_a^-$, respectively. The required operator is
\begin{equation}
 \hat f_a=a_a+b_a\hat{\mathbf s}\hat{\mathbf i}_a,
 \tag{151.4}
\end{equation}
where $\hat{\mathbf i}_a$ and $\hat{\mathbf s}$ are the nuclear- and neutron-spin operators, and the coefficients $a_a$ and $b_a$ are
\begin{equation}
 a_a=\frac1{2i_a+1}\big[(i_a+1)f_a^++i_af_a^-\big],\qquad
 b_a=\frac2{2i_a+1}(f_a^+-f_a^-).
 \tag{151.5}
\end{equation}

\SourcePage{767}{770}
This is readily verified by noting that, for specified $j$, the eigenvalue of $\mathbf s\mathbf i$ is
\[
 \mathbf s\mathbf i=\frac12\left[j(j+1)-i(i+1)-\frac34\right].
\]

The operators (151.4), with their matrix elements for the transition under consideration, must be substituted for the $f_a$ in (151.3). If the incident neutrons and the target nuclei are unpolarized, the scattering cross section must be averaged accordingly.

\subsection*{Problems}

\ProblemHeading{1.} Average formula (151.3), assuming that the spin directions of the neutrons and nuclei are distributed completely at random. All nuclei in the molecule are different.

\SolutionHeading The averages over the neutron- and nuclear-spin directions are independent, and each individual spin averages to zero; hence $\overline{\mathbf s\mathbf i_a}=0$. If the molecule contains no identical atoms, there is no exchange interaction between the nuclear spins. Because their direct interaction is negligible, the spin directions of different nuclei in the molecule may be regarded as independent; products of the form $(\mathbf s\mathbf i_1)(\mathbf s\mathbf i_2)$ therefore also vanish on averaging. For the squares $(\mathbf s\mathbf i)^2$, however,
\[
 \overline{(\mathbf s\mathbf i)^2}=\frac13\mathbf s^2\mathbf i^2
 =\frac{s(s+1)i(i+1)}3=\frac{i(i+1)}4.
\]
The averaged cross section is consequently
\[
 d\sigma_n=M_\text{mol}^2\frac{p'}p\left\{
 \left|\sum_a\frac{a_a}{M_a}\langle n|e^{-i\mathbf q\mathbf R_a}|0\rangle\right|^2
 +\frac14\sum_a\frac{i_a(i_a+1)}{M_a^2}b_a^2
 \left|\langle n|e^{-i\mathbf q\mathbf R_a}|0\rangle\right|^2\right\}do.
\]

\ProblemHeading{2.} Apply formula (151.3) to the scattering of slow neutrons by para- and orthohydrogen (J.~Schwinger, E.~Teller, 1937).

\SolutionHeading Before the matrix elements of the spin operators are evaluated, expression (151.3) for scattering by an $\mathrm H_2$ molecule is
\begin{equation}
 d\sigma_n=\frac{16p'}{9p}\left|a\langle n|e^{-i\mathbf q\mathbf r/2}+e^{i\mathbf q\mathbf r/2}|0\rangle
 +b\hat{\mathbf s}\langle n|\hat{\mathbf i}_1e^{-i\mathbf q\mathbf r/2}+\hat{\mathbf i}_2e^{i\mathbf q\mathbf r/2}|0\rangle\right|^2do,
 \tag{1}
\end{equation}
\[
 a=\frac14(3f^++f^-),\qquad b=f^+-f^-
\]
where $\pm\mathbf r/2$ are the radius vectors of the two nuclei relative to the molecule's center of mass.

The rotational and vibrational states of the molecule are specified by the quantum numbers $K$, $M_K$, and $v$, whose collection is denoted by $n$ in (1). In the ground electronic state of $\mathrm H_2$, even $K$ values are possible only for total nuclear spin $I=0$, para-hydrogen, while odd $K$ values require $I=1$, ortho-hydrogen (see §~86). Two cases must therefore be distinguished: (1) transitions between rotational states whose $K$ values have the same parity, possible only without a change of $I$ (ortho--ortho and para--para transitions); and (2) transitions between states whose $K$ values have different parity, possible only with a change of $I$ (ortho--para and para--ortho transitions).

\SourcePage{768}{771}
In the first case,
\[
 \langle n|e^{-i\mathbf q\mathbf r/2}|0\rangle
 =\langle n|e^{i\mathbf q\mathbf r/2}|0\rangle
 =\left\langle n\left|\cos\frac{\mathbf q\mathbf r}{2}\right|0\right\rangle
\]
(recall that the rotational wave function is multiplied by $(-1)^K$ when the sign of $\mathbf r$ is reversed). The spin operator in (1) then becomes $2a+b\hat{\mathbf s}\hat{\mathbf I}$, where $\hat{\mathbf I}=\hat{\mathbf i}_1+\hat{\mathbf i}_2$. This operator is diagonal in $I$, in accordance with the statement above. Averaging the square $(2a+b\mathbf s\mathbf I)^2$ as in Problem~1 gives
\[
 4a^2+\frac{b^2}{4}I(I+1).
\]
The result is
\begin{equation}
 d\sigma_n=\frac{4p'}{9p}\left|\left\langle n\left|\cos\frac{\mathbf q\mathbf r}{2}\right|0\right\rangle\right|^2
 \big\{(3f^++f^-)^2+I(I+1)(f^+-f^-)^2\big\}do.
 \tag{2}
\end{equation}

In the second case,
\[
 \langle n|e^{i\mathbf q\mathbf r/2}|0\rangle
 =-\langle n|e^{-i\mathbf q\mathbf r/2}|0\rangle
 =i\left\langle n\left|\sin\frac{\mathbf q\mathbf r}{2}\right|0\right\rangle,
\]
and the spin operator in (1) reduces to $\hat{\mathbf s}(\hat{\mathbf i}_1-\hat{\mathbf i}_2)$; its matrix elements are exclusively off-diagonal in $I$. The squared modulus of these elements, summed over every possible value of the final-state total-spin projection $I'$, is calculated as the mean value, or diagonal element, of the square $(\mathbf s,\mathbf i_1-\mathbf i_2)^2$ (see the note on p.~704) and equals
\[
 \overline{(\mathbf s,\mathbf i_1-\mathbf i_2)^2}
 =\frac13\cdot\frac34\,\overline{(\mathbf i_1-\mathbf i_2)^2}
 =\frac14(2\mathbf i_1^2+2\mathbf i_2^2-\mathbf I^2)
 =\frac14[3-I(I+1)].
\]
It follows that
\begin{equation}
 d\sigma_n=(1)(3)\frac{4p'}{9p}
 \left|\left\langle n\left|\sin\frac{\mathbf q\mathbf r}{2}\right|0\right\rangle\right|^2(f^+-f^-)^2do,
 \tag{3}
\end{equation}
where the coefficient (1) applies to ortho--para transitions and the coefficient (3) to para--ortho transitions.

If the neutrons are so slow that their wavelength is also large compared with the molecular dimensions, one may set $\cos(qr/2)=1$ and $\sin(qr/2)=0$ in the matrix elements of (2) and (3). They all then vanish except for the diagonal element 00; naturally, only elastic scattering is possible under these conditions. Its cross section is
\[
 d\sigma_e=\frac49\big[(3f^++f^-)^2+I(I+1)(f^+-f^-)^2\big]do.
\]

\ProblemHeading{3.} Determine the neutron-scattering cross section for a bound proton treated as an isotropic three-dimensional oscillator of frequency $\omega$ (E.~Fermi, 1936).

\SolutionHeading Treating the proton as vibrating about a point fixed in space, we must set $M_\text{mol}=M$ and $M_a=M/2$ in (151.3), in accordance with its derivation; $M$ is the proton mass. Then
\[
 d\sigma_n=\frac{p'}p\frac{\sigma_0}{\pi}\sum
 \left|\int e^{-i\mathbf q\mathbf r}\psi_{000}(\mathbf r)\psi_{n_1n_2n_3}(\mathbf r)\,dV\right|^2do,
\]

\SourcePage{769}{772}
where $\sigma_0=4\pi|f|^2$ is the free-proton scattering cross section, while $\psi_{n_1n_2n_3}$ are eigenfunctions of the three-dimensional oscillator corresponding to the energy levels $E_n=\hbar\omega(n+3/2)$. The sum extends over all $n_1,n_2,n_3$ with specified sum $n_1+n_2+n_3=n$. The functions $\psi_{n_1n_2n_3}$ are products of wave functions of three linear oscillators (see Problem~4, §~33). The required integral therefore separates into a product of three integrals of the form
\[
 \int_{-\infty}^{\infty}\exp\left(-\frac{iq_xx}{2}-\frac{\alpha^2x^2}{2}-\frac{\alpha^2x^2}{2}\right)H_{n_1}(\alpha x)\,dx
\]
where $\alpha=\sqrt{M\omega/\hbar}$. They are evaluated by substituting the representation (a.4) for $H_n(x)$ and integrating by parts $n$ times. This gives
\[
 d\sigma_n=\frac1\pi\frac{v'}v\frac{\sigma_0}{2^n\alpha^{2n}}
 \sum\frac{q_x^{2n_1}q_y^{2n_2}q_z^{2n_3}}{n_1!n_2!n_3!}
 \exp\left(-\frac{q^2}{2\alpha^2}\right)do.
\]
The sum is evaluated with the multinomial formula, and finally
\[
 d\sigma_n=\frac{\sigma_0}{\pi n!}\sqrt{\frac{E'}E}
 \left(\frac{q^2}{2\alpha^2}\right)^n
 \exp\left(-\frac{q^2}{2\alpha^2}\right)do.
\]
In particular, the elastic-scattering cross section, for $n=0$ and $E=E'$, is
\[
 d\sigma_e=\frac{\sigma_0}{\pi}\exp\left(-\frac{q^2}{2\alpha^2}\right)do,
 \qquad
 \sigma_e=\sigma_0\frac{\hbar\omega}{E}
 \left[1-\exp\left(-\frac{4E}{\hbar\omega}\right)\right].
\]
If $E/\hbar\omega\longrightarrow0$, then $\sigma_e\longrightarrow4\sigma_0$.
